\documentclass[11pt]{article}

% ============================================================================
% DIMENSIONI PERSONALIZZATE (Modifica qui!)
% ============================================================================
% Esempio per 200 pagine su carta bianca (Dorso = 0.45in)
% Larghezza totale = 6 (front) + 6 (back) + 0.45 (spine) + 0.25 (bleed) = 12.7in
\usepackage[paperwidth=12.7in, paperheight=9.25in, margin=0in]{geometry}

% ============================================================================
% PACKAGES
% ============================================================================
\usepackage[utf8]{inputenc}
\usepackage[T1]{fontenc}
\usepackage{graphicx}
\usepackage{xcolor}
\usepackage{mathpazo} % Mantiene il font Palatino sobrio di AstDyn
\usepackage{tikz}

% Colore tema (usiamo il blu link di AstDyn o nero sobrio)
\definecolor{themecolor}{rgb}{0.0,0.1,0.4}

% Modifica questo valore con lo spessore reale del tuo dorso
\newcommand{\spinewidth}{0.45in} 

\begin{document}
\thispagestyle{empty}

\begin{tikzpicture}[remember picture, overlay]
    % Sfondo totale (Bianco o crema sobrio)
    \fill[white] (current page.north west) rectangle (current page.south east);

    % ------------------------------------------------------------------------
    % QUARTA DI COPERTINA (Back Cover - Sinistra)
    % ------------------------------------------------------------------------
    \node[anchor=north west, minimum width=6in, minimum height=9in, text width=5in, align=left, inner sep=1in] 
        at (current page.north west) {
            {\Large\bfseries\color{themecolor} About AstDyn}\\[1ex]
            \small
            This manual provides a comprehensive scientific reference for the \textbf{AstDyn} library, 
            part of the ITALOccult framework. It covers the theoretical foundations of celestial mechanics, 
            orbit determination algorithms, numerical integration methodologies, and practical implementation details.
            Designed for researchers, engineers, and astronomy enthusiasts, this guide ensures high-precision dynamics 
            for asteroid and comet simulations.
            
            \vspace{2cm}
            % Spazio per il codice a barre di Amazon
            \hfill \fbox{\phantom{\rule{2in}{1.2in}}}
        };

    % ------------------------------------------------------------------------
    % DORSO (Spine - Centro)
    % ------------------------------------------------------------------------
    % Linea di demarcazione del dorso (opzionale, utile per verifica)
    \draw[gray!30, dashed] ({6in + 0.125in}, 0) -- ({6in + 0.125in}, -9.25in);
    \draw[gray!30, dashed] ({6in + 0.125in + \spinewidth}, 0) -- ({6in + 0.125in + \spinewidth}, -9.25in);

    \node[anchor=center, minimum width=\spinewidth, minimum height=9in, rotate=-90] 
        at ({6in + 0.125in + \spinewidth/2}, -9in/2 - 0.125in) {
            {\small\fboxsep=2pt\color{themecolor}\bfseries AstDyn: Scientific Reference Manual}
            \hfill {\small Michele Bigi}
        };

    % ------------------------------------------------------------------------
    % PRIMA DI COPERTINA (Front Cover - Destra)
    % ------------------------------------------------------------------------
    \node[anchor=north east, minimum width=6in, minimum height=9in, inner sep=0pt] 
        at ({12in + 0.125in}, -0.125in) {
            \begin{tikzpicture}[remember picture, overlay]
                % Fascia decorativa (Blu notte)
                \fill[themecolor] (0,0) rectangle (6in, -2.5in);
                
                % Titolo Principale
                \node[anchor=west, text width=5in, font=\fontsize{36}{40}\selectfont\bfseries\color{white}] 
                    at (0.5in, -1.25in) {AstDyn};
                
                % Sottotitolo
                \node[anchor=west, text width=5in, font=\fontsize{14}{16}\selectfont\color{white!90}] 
                    at (0.5in, -2in) {Scientific Reference Manual};

                % Sottotitolo esteso
                \node[anchor=west, text width=5in, font=\large\scshape\color{themecolor}] 
                    at (0.5in, -3in) {The ITALOccult Framework for\\ High-Precision Asteroid Dynamics};

                % Grafico o Diagramma (Se ne hai uno PNG/JPG di AstDyn, inseriscilo qui)
                \node[anchor=center] at (3in, -6in) {
                    % \includegraphics[width=4in]{tuo_grafico_orbita.png} 
                    % Rimuovi il commento sopra se hai un'immagine
                    \fbox{\phantom{\rule{4in}{3in}}} % Segnaposto
                };

                % Autore
                \node[anchor=south east, font=\large\bfseries\color{themecolor}] 
                    at (5.5in, -8.5in) {Michele Bigi};
            \end{tikzpicture}
        };

\end{tikzpicture}

\end{document}
